%%%%%%%%%%%%%%%%%%%%%%%%%%%%%%%%%%%%%%%%%%%%%%%%%%%%%%%%%%%%%%%%%%%%%%%%%%%%%%%%
\documentclass[11pt]{article} % Dokumentenklasse

\usepackage[utf8]{inputenc} % Textkodierung: UTF-8
\usepackage[T1]{fontenc} % Zeichensatzkodierung

\usepackage[USenglish]{babel}% http://ctan.org/pkg/babel
\usepackage{graphicx} % Grafiken
\usepackage[absolute]{textpos} % Positionierung

% Schriftart Helvetica:
\usepackage[scaled]{helvet}
\renewcommand{\familydefault}{\sfdefault}

\usepackage{calc} % Berechnungen
\usepackage{tabto} % Tabulatoren
\usepackage{parskip}

\usepackage{enumitem}

% Debugging:
%\usepackage{layout} % Layout-Informationen
%\usepackage{printlen} % Längenwerte ausgeben

%%%%%%%%%%%%%%%%%%%%%%%%%%%%%%%%%%%%%%%%%%%%%%%%%%%%%%%%%%%%%%%%%%%%%%%%%%%%%%%%
% EINSTELLUNGEN
%%%%%%%%%%%%%%%%%%%%%%%%%%%%%%%%%%%%%%%%%%%%%%%%%%%%%%%%%%%%%%%%%%%%%%%%%%%%%%%%

% Seitenränder:
\newcommand{\SeitenrandOben}{43.5mm}
\newcommand{\SeitenrandRechts}{20mm}
\newcommand{\SeitenrandLinks}{20mm}
\newcommand{\SeitenrandUnten}{10mm}

\newcommand{\UniversitaetLogoBreite}{19mm}
\newcommand{\UniversitaetLogoHoehe}{1cm}

\usepackage[a4paper,
top=\SeitenrandOben,
bottom=\SeitenrandUnten,
inner=\SeitenrandLinks,
outer=\SeitenrandRechts,
foot=0cm,
head=0cm
]{geometry}

\textblockorigin{\SeitenrandLinks}{\SeitenrandOben} % Ursprung für Positionierung

\setlength{\parindent}{0pt}
%\setlength{\baselineskip}{32pt}
\setlength{\parskip}{\baselineskip}
\TabPositions{4cm}
\pagestyle{empty}

%%%%%%%%%%%%%%%%%%%%%%%%%%%%%%%%%%%%%%%%%%%%%%%%%%%%%%%%%%%%%%%%%%%%%%%%%%%%%%%%
% General stuff
%%%%%%%%%%%%%%%%%%%%%%%%%%%%%%%%%%%%%%%%%%%%%%%%%%%%%%%%%%%%%%%%%%%%%%%%%%%%%%%%
\newcommand{\Problem}[1]{\paragraph*{Problem #1:}\qquad}
\newcommand{\Topic}[1]{
	\newpage
	\section*{#1}}

\newcommand{\Given}{\textbf{Given:\qquad\qquad}}
\newcommand{\Searched}{\textbf{Searched:\qquad}}
\newcommand{\Solution}{\textbf{Solution:\qquad}}

%%%%%%%%%%%%%%%%%%%%%%%%%%%%%%%%%%%%%%%%%%%%%%%%%%%%%%%%%%%%%%%%%%%%%%%%%%%%%%%%
% Math stuff
%%%%%%%%%%%%%%%%%%%%%%%%%%%%%%%%%%%%%%%%%%%%%%%%%%%%%%%%%%%%%%%%%%%%%%%%%%%%%%%%
\usepackage{amsmath}
\usepackage{amssymb}

\newcommand{\R}{\mathbb{R}}
\newcommand{\Vector}[1]{\R^{#1}}
\newcommand{\Matrix}[2]{\R^{#1 \times #2}} % !!! DON'T TOUCH !!!
%%%%%%%%%%%%%%%%%%%%%%%%%%%%%%%%%%%%%%%%%%%%%%%%%%%%%%%%%%%%%%%%%%%%%%%%%%%%%%%%


\newcommand{\ExerciseNumber}{08}

\newcommand{\PersonOne}{Marcel Bruckner (03674122)}
\newcommand{\PersonTwo}{Julian Hohenadel (03673879)}
\newcommand{\PersonThree}{Kevin Bein (03707775)}


%%%%%%%%%%%%%%%%%%%%%%%%%%%%%%%%%%%%%%%%%%%%%%%%%%%%%%%%%%%%%%%%%%%%%%%%%%%%%%%%
% DOKUMENT
%%%%%%%%%%%%%%%%%%%%%%%%%%%%%%%%%%%%%%%%%%%%%%%%%%%%%%%%%%%%%%%%%%%%%%%%%%%%%%%%

\begin{document}

%%%%%%%%%%%%%%%%%%%%%%%%%%%%%%%%%%%%%%%%%%%%%%%%%%%%%%%%%%%%%%%%%%%%%%%%%%%%%%%%
\begin{textblock*}{\UniversitaetLogoBreite}[1,0](\textwidth-1mm, 2cm-\SeitenrandOben)%
	\raggedleft\includegraphics{../Ressources/Universitaet_Logo_RGB.pdf}%
\end{textblock*}


\begin{textblock*}{\textwidth}[0,0](0cm, 0cm)%
	{\fontsize{24pt}{26pt}\selectfont\textbf{Exercise}}
	
	\vspace*{14pt}
	{\fontsize{18pt}{27pt}\selectfont\textbf{\ExerciseNumber}}
\end{textblock*}

\vspace*{92.2mm}
\fontsize{15pt}{17.5pt}\selectfont%
TUM Department of Informatics

\renewcommand{\baselinestretch}{1.47}
\normalsize\selectfont
\vspace*{17.1mm}
\textbf{Supervised by}\tab
\begin{minipage}[t]{\textwidth-\CurrentLineWidth}
	Prof. Dr. Stephan Günnemann\\
	Informatics 3 - Professorship of Data Mining and Analytics\strut
\end{minipage}

\vspace*{4.3mm}
\textbf{Submitted by}\tab
\begin{minipage}[t]{\textwidth-\CurrentLineWidth}
	\PersonOne\\
	\PersonTwo\\
	\PersonThree
\end{minipage}

\vspace*{-1mm}
\textbf{Submission date}\tab 
\begin{minipage}[t]{\textwidth-\CurrentLineWidth}
	Munich, \today
\end{minipage}
\newpage % !!! DON'T TOUCH !!!
%%%%%%%%%%%%%%%%%%%%%%%%%%%%%%%%%%%%%%%%%%%%%%%%%%%%%%%%%%%%%%%%%%%%%%%%%%%%%%%%

%%%%%%%%%%%%%%%%%%%%%%%%%%%%%%%%%%%%%%%%%%%%%%%%%%%%%%%%%%%%%%%%%%%%%%%%%%%%%%%%
% !!! HOMEWORK STARTS HERE !!!
%%%%%%%%%%%%%%%%%%%%%%%%%%%%%%%%%%%%%%%%%%%%%%%%%%%%%%%%%%%%%%%%%%%%%%%%%%%%%%%%
%
\Topic{SVM and Kernels}
%
\Problem{1}
%
\begin{flushleft}
Similarities: \\
Both try to find a fitting hyperplane which seperates the data classes.\\
Difference: \\
SVM tries to maximize the margin from the hyperplane to the data points, perceptron algorithms only care about a valid seperation of the data classes.
\end{flushleft}
%
%
\Problem{2}
%
\begin{flushleft}
a)\\
$g(\alpha)$ vectorized definition:
\begin{align*}
g(\alpha) = \frac{1}{2}\alpha^{T}Q\alpha + \alpha^{T}1_{N}
\end{align*}
$g(\alpha)$ standard definition:
\begin{align*}
g(\alpha) = \sum_{i=1}^{N}\alpha_{i}-\frac{1}{2}\sum_{i=1}^{N}\sum_{j=1}^{N}
y_{i}y_{j}\alpha_{i}\alpha_{j}x^{T}_{i}x_{j}
\end{align*}
$y$ is a vector of dimension $N\times 1$\\
$x$ is a matrix of dimension $N\times M$\\
$\sum_{i=1}^{N}\sum_{j=1}^{N}y_{i}y_{j}$ is equivalent to $yy^{T}$ 
(dimension is $N\times N$)\\
$\sum_{i=1}^{N}\sum_{j=1}^{N}x^{T}_{i}x_{j}$ is equivalent to $XX^{T}$ 
(dimension is $N\times N$)\\
$\sum_{i=1}^{N}\sum_{j=1}^{N}y_{i}y_{j}x^{T}_{i}x_{j}$ is the Hadamard product so:
$[yy^{T} \odot XX^{T}]$\\
Take the $-1$ scalar from the standard definition into the matrix: 
$[-yy^{T} \odot XX^{T}] = Q$\\
$\implies \frac{1}{2}\alpha^{T}Q\alpha \equiv \frac{1}{2}\alpha^{T}[-yy^{T} 
\odot XX^{T}]\alpha \equiv -\frac{1}{2}\sum_{i=1}^{N}\sum_{j=1}^{N}y_{i}y_{j}\alpha_{i}
\alpha_{j}x^{T}_{i}x_{j}$\\
$\alpha^{T}1_{N} \equiv \sum_{i=1}^{N}\alpha_{i}$ is trivial.\\
$\implies$ $g(\alpha)$ vectorized definition $\equiv$ $g(\alpha)$ standard definition
\end{flushleft}
\begin{flushleft}
b)\\
If the search for a local maximizer of $g$ returns the global maximum of $g$, that
means that the maximization problem is concave.\\
To proove this claim $Q$ needs to be negativ (semi)definite (NSD).\\
For $Q$ to be NSD: $\forall \alpha \in \R^{N}: \alpha^{T}Q\alpha \leq 0$ 
needs to hold.\\
\begin{align}
\alpha^{T}Q\alpha &\leq 0\\
-\sum_{i=1}^{N}\sum_{j=1}^{N}y_{i}y_{j}\alpha_{i}\alpha_{j}x^{T}_{i}x_{j} &\leq 0\\
-\sum_{i=1}^{N}\sum_{j=1}^{N}(y_{i}\alpha_{i}x_{i})^{T}(y_{j}\alpha_{j}x_{j}) &\leq 0\\
-(y\odot\alpha)^{T}XX^{T}(y\odot\alpha) &\leq 0\\
-(y\odot\alpha)^{T}(X^{T})^{T}X^{T}(y\odot\alpha) &\leq 0\\
-(X^{T}(y\odot\alpha))^{T}(X^{T}(y\odot\alpha)) &\leq 0\\
-(X^{T}(y\odot\alpha))^{2} &\leq 0\\
-(\geq 0) &\leq 0 \qquad \square
\end{align}
(3): $y_{i}$ and $\alpha_{i}$ can be dragged inside here because they are scalars.\\
(4): The whole expression returns a scalar, that's why reshaping is done this way:\\
$dim((y\odot\alpha)^{T})= 1 \times N$\\
$dim(X)= N \times 1$\\
$dim(X^{T})= 1 \times N$\\
$dim((y\odot\alpha))= N \times 1$\\
(5): $X=(X^{T})^{T}$\\
(6): $(AB)^{T}=B^{T}A^{T}$\\
(7): $(...)^{2} \geq 0 $, (as long as $...$ is not complex)\\
(8): Proofs that $Q$ is NSD.\\
$Q$ is NSD, that means the maximization problem is in fact concave so local maxima = global maxima.
\end{flushleft}
%
%
\Problem{3}
%
%
%
\Problem{4}
%
%
%
\Problem{5}
%
\begin{flushleft}
"A kernel is valid if it corresponds to an inner product in some feature space."\\
$x_{1}^{T}x_{2}$ is a valid kernel because it is a scalar product of the input vectors.\\
$a_{0}$ is a constant term and a valid kernel because it can be generated by:\\
$k(x_i,x_j) = \phi(x_i)^T\phi(x_j)$ with $\phi(x) = \sqrt{a_0} \implies k(x_i,x_j) = a_0$\\
The same holds for $a_i$ so $a_i(x_i^T x_j)^i + a_0$ can be represented as kernels.\\
Sums and multiplications of kernels are kernel preserving (also scalars are $\geq 0$)\\
$\implies$ Gram matrix is still PSD.
\end{flushleft}
%
%
\Problem{6}
%
\begin{flushleft}
The "trick" for this Problem:\\
"The Maclaurin series for $\frac{1}{1-x}$ is the geometric series:
$1 + x + x^2 + x^3 + x^4 ...$"
\begin{align*}
k(x_1,x_2)&=\frac{1}{1-x_1 x_2}\\
&= \sum_{i=0}^{\infty}x_1^i x_2^i\\
&= \phi(x_1)^T \phi(x_2)\\
\text{with } \phi(x)&=(1,x,x^2,x^3,x^4, ...)
\end{align*}
\end{flushleft}
%
%
%%%%%%%%%%%%%%%%%%%%%%%%%%%%%%%%%%%%%%%%%%%%%%%%%%%%%%%%%%%%%%%%%%%%%%%%%%%%%%%%
% !!! HOMEWORK ENDS HERE !!!
%%%%%%%%%%%%%%%%%%%%%%%%%%%%%%%%%%%%%%%%%%%%%%%%%%%%%%%%%%%%%%%%%%%%%%%%%%%%%%%%

%%%%%%%%%%%%%%%%%%%%%%%%%%%%%%%%%%%%%%%%%%%%%%%%%%%%%%%%%%%%%%%%%%%%%%%%%%%%%%%%
\newpage

\vspace*{-15.8mm}
\fontsize{19pt}{21pt}\selectfont

\vspace{25.3mm}
Appendix

\normalsize\selectfont
\vspace{13.2mm}
We confirm that the submitted solution is original work and was written by us without further assistance. Appropriate credit has been given where reference has been made to the work of others.

\vspace{18.1mm}
\rule[-3.7mm]{\linewidth}{0.5pt}
Munich, \today, Signature \PersonOne

\vspace{18.1mm}
\rule[-3.7mm]{\linewidth}{0.5pt}
Munich, \today, Signature \PersonTwo

\vspace{18.1mm}
\rule[-3.7mm]{\linewidth}{0.5pt}
Munich, \today, Signature \PersonThree
 % !!! DON'T TOUCH !!!
%%%%%%%%%%%%%%%%%%%%%%%%%%%%%%%%%%%%%%%%%%%%%%%%%%%%%%%%%%%%%%%%%%%%%%%%%%%%%%%%

\end{document}
